% 恒星光谱（综述）
% license CCBYSA3
% type Wiki

本文根据 CC-BY-SA 协议转载翻译自维基百科\href{https://en.wikipedia.org/wiki/Stellar_classification}{相关文章}。

在天文学中，恒星分类是根据恒星的光谱特性对恒星进行的分类。来自恒星的电磁辐射通过棱镜或衍射光栅被分解成光谱，呈现出类似彩虹的连续颜色，其间夹杂着光谱线。每一条光谱线都对应一种特定的化学元素或分子，而谱线的强度则反映该元素的丰度。不同光谱线强度的变化主要由恒星光球层的温度决定，尽管在某些情况下也确实存在真实的元素丰度差异。恒星的光谱型是一种简短的代码，主要概括了其电离状态，从而为光球层温度提供一个客观的度量。

目前，大多数恒星采用摩根–基南（Morgan–Keenan，MK）分类系统进行分类，该系统使用字母 O、B、A、F、G、K、M，按从最热（O 型）到最冷（M 型）的顺序排列。每一个字母类别又进一步用一个数字细分，其中 0 表示最热，9 表示最冷（例如 A8、A9、F0、F1 就构成了一个由热到冷的序列）。此外，该序列还扩展出 W、S、C 三个类别，用于那些不符合经典体系的恒星。一些恒星遗迹或质量明显偏离恒星范围的天体也被赋予了字母标记：D 表示白矮星，L、T、Y 表示棕矮星（以及部分系外行星）。

在 MK 系统中，还会在光谱型后加上一个用罗马数字表示的光度级。光度级依据恒星光谱中某些吸收线的宽度来判定，这些线宽随恒星大气密度变化，因此可用于区分巨星与矮星。其中，光度级 0 或 Ia+ 表示特超巨星，I 表示超巨星，II 表示亮巨星，III 表示普通巨星，IV 表示亚巨星，V 表示主序星，sd（或 VI）表示亚矮星，而 D（或 VII）表示白矮星。太阳的完整光谱分类为 G2V，这表明太阳是一颗主序星，其表面温度约为 5800 K。
\subsection{传统颜色描述}
传统的恒星颜色描述只考虑恒星光谱的峰值位置。然而实际上，恒星会在整个电磁波谱范围内辐射。由于所有光谱颜色叠加后呈现为白色，人眼实际看到的恒星颜色要比传统颜色描述所暗示的浅得多。这种“明亮度”的特性说明，将恒星简单地赋予某种光谱颜色可能具有误导性。在排除暗光条件下的颜色对比效应后，在一般观测条件中，并不存在真正的绿色、青色、靛色或紫色恒星。像太阳这样的“黄色”矮星实际上呈现为白色；“红矮星”呈现为较深的黄色或橙色；而“棕矮星”并不会真的呈现棕色，而是假想中在近距离观察时可能显得暗红色，或呈灰色/黑色。
\subsection{现代分类}
现代恒星分类体系被称为摩根–基南（Morgan–Keenan，MK）分类系统。在该系统中，每一颗恒星都会被赋予一个光谱型（源自较早的哈佛光谱分类体系，该体系本身并不包含光度级$^{[1]}$），并再配以一个用罗马数字表示的光度级（如下所述），二者共同构成恒星的完整光谱类型。

其他现代恒星分类体系（例如 UBV 系统）则基于颜色指数，即通过测量两个或多个波段星等之间的差值来进行分类$^{[2]}$。这些数值通常以诸如 U−V、B−V这样的形式给出，表示通过两种标准滤光片（例如紫外、蓝光和可见光）所测得的颜色差异。
\subsubsection{哈佛光谱分类}
哈佛体系是一种一维的恒星分类方案，由天文学家安妮·江普·坎农（Annie Jump Cannon）提出，她对德雷珀（Draper）早期按字母顺序排列的分类体系进行了重新排序与简化（见“历史”部分）。在该体系中，恒星根据其光谱特征被划分为若干组，每一组用一个字母表示，并可辅以数字细分。

主序星的表面温度大致分布在 2000–50 000 K 之间，而演化程度更高的恒星——尤其是新近形成的白矮星——其表面温度则可超过 100 000 K$^{[3]}$。从物理意义上看，这些光谱类别反映的是恒星大气的温度，并通常按由热到冷的顺序排列。
\begin{figure}[ht]
\centering
\includegraphics[width=14.25cm]{./figures/1615d7b930f17d79.png}
\caption{} \label{fig_Stcl_1}
\end{figure}
用于记忆光谱型字母由热到冷排列顺序的传统助记口诀是“Oh, Be A Fine Guy/Girl: Kiss Me!”（“哦，做个好小伙/好姑娘：吻我吧！”）$^{[12]}$。
尽管在天文学课程和相关组织举办的活动中，人们提出过许多替代性的助记口诀，但这一传统口诀至今仍然是最为流行的$^{[13][14]}$。

光谱类型 O 到 M（以及后文讨论的其他更为专门的类型）都会再用阿拉伯数字 0–9进行细分，其中 0 表示该类型中最热的恒星。例如，A0 表示 A 型中最热的恒星，而 A9 则表示最冷的 A 型恒星。允许使用小数；例如，恒星 Mu Normae 的光谱型被定为 O9.7 $^{[15]}$。太阳的光谱型为 G2$^{[16]}$。

尽管哈佛光谱分类实际上反映的是恒星的表面温度或光球温度（更准确地说是其有效温度），但这一点在该体系建立之初并未被完全理解。不过，到第一张赫罗图（Hertzsprung–Russell diagram）于 1914 年前后提出时，人们已经普遍怀疑这一对应关系是成立的$^{[17]}$。在 1920 年代，印度物理学家梅格纳德·萨哈（Meghnad Saha）在物理化学中分子解离理论的基础上，发展出了电离理论并将其推广到原子的电离问题。他首先将该理论应用于太阳色球层，随后又应用于恒星光谱$^{[18]}$。

随后，哈佛天文学家塞西莉亚·佩恩（Cecilia Payne）证明O–B–A–F–G–K–M这一光谱序列本质上就是一个温度序列$^{[19]}$。由于该分类序列的建立早于人们对其“温度本质”的认识，将一条光谱归入某一具体亚型（如 B3 或 A7）在很大程度上依赖于对恒星光谱中吸收特征强度的（带有一定主观性的）估计。因此，这些亚型并不是在数学意义上等间隔划分的。
\subsubsection{摩根–基南（Morgan–Keenan）分类}
耶基斯光谱分类，也称为 MK 分类，或 Morgan–Keenan（亦称 MKK，Morgan–Keenan–Kellman）系统$^{[20][21]}$，是一种由威廉·威尔逊·摩根（William Wilson Morgan）、菲利普·C·基南（Philip C. Keenan）和伊迪丝·凯尔曼（Edith Kellman）于 1943 年在耶基斯天文台提出的恒星光谱分类体系$^{[22]}$。这是一个二维分类体系（温度与光度），其依据是对恒星温度和表面重力敏感的光谱线；而表面重力又与恒星的光度密切相关（相比之下，哈佛分类仅基于表面温度）。在 1953 年对标准星表和分类标准进行若干修订之后，该体系被正式命名为摩根–基南分类（MK）$^{[23]}$，并一直沿用至今。

表面重力更高、密度更大的恒星，其光谱线会表现出更显著的压力展宽。由于巨星的半径远大于质量相近的矮星，其表面重力（以及由此产生的压力）要低得多。因此，光谱中的差异可以被解释为光度效应，从而仅通过分析光谱本身就能够确定恒星的光度级。

目前区分出了若干不同的光度级，具体列于下表中$^{[24]}$。
\begin{figure}[ht]
\centering
\includegraphics[width=14.25cm]{./figures/73151fcf04a39c9e.png}
\caption{} \label{fig_Stcl_2}
\end{figure}
允许存在边缘情况；例如，一颗恒星可能既可被归为超巨星也可被归为亮巨星，或者介于亚巨星与主序星之间。在这些情况下，会在两个光度级之间使用两种特殊符号：
\begin{itemize}
\item 斜杠（/）：表示恒星可能属于其中任意一个光度级。
\item 连字符（-）：表示恒星处在两个光度级之间的过渡状态。
\end{itemize}
例如，被标记为 A3-4III/IV 的恒星，表示其光谱型介于 A3 与 A4 之间，同时其光度级介于巨星（III）与亚巨星（IV）之间，或可能属于其中之一。

此外，还曾使用亚矮星（sub-dwarf）光度级：

VI 表示亚矮星（其光度略低于主序星）。

名义上的光度级 VII（有时甚至更高的罗马数字）如今已很少用于白矮星或“热亚矮星”的分类，因为主序星和巨星所采用的温度字母体系已不再适用于白矮星。

在某些情况下，字母 a 和 b 也会用于非超巨星的光度级中；例如，一颗光度略低于典型巨星的恒星可被标为 IIIb，而 IIIa 则表示其光度略高于典型巨星$^{[34]}$。

一小部分 V 型极端恒星在 He II λ4686 吸收谱线处表现出很强的吸收特征，因此被赋予 Vz 的分类标记。一个典型例子是恒星 HD 93129 B$^{[35]}$。
\subsubsection{光谱异常}
还可以在光谱型之后附加小写字母形式的补充标记，用以指出光谱中的某些特殊或异常特征$^{[36]}$。
\begin{table}[ht]
\centering
\caption{请输入表格标题}\label{tab_Stcl1}
\begin{tabular}{|c|c|}
\hline
\textbf{代码} & \textbf{恒星的光谱异常} \\
\hline
: & 不确定的光谱值$^{[24]}$\\
\hline
... & 存在未描述的光谱异常\\
\hline
! & 特殊异常\\
\hline
comp & 复合光谱[37]\\
\hline
e & 存在发射线[37]\\
\hline
[e] & 存在“禁戒”发射线\\
\hline
er & “反向”发射线中心比边缘弱\\
\hline
eq & 具有P Cygni轮廓的发射线\\
\hline
f & N III 和 He II 发射$^{[24]}$\\
\hline
f & NIV 4058Å线比NIII 4634Å、4640Å和4642Å线更强$^{[38]}$\\
\hline
f+ & 除了N III线外，还发射Si IV 4089Å和4116Å线$^{[38]}$\\
\hline
f? & C III 4647–4650–4652Å发射线的强度与N III线相当$^{[39]}$\\
\hline
(f) & N III发射，He II缺失或吸收较弱\\
\hline
(f+) & $^{[40]}$\\
\hline
((f)) & 显示强烈的He II吸收，伴随微弱的N III发射$^{[41]}$\\
\hline
((f*)) & $^{[40]}$\\
\hline
h & WR恒星以氢发射线为特征。$^{[42]}$\\
\hline
ha & WR恒星的氢线在吸收和发射中均可见。$^{[42]}$\\
\hline
He wk & 弱氦线\\
\hline
k & 带有星际吸收特征的光谱\\
\hline
m & 增强的金属特征$^{[37]}$\\
\hline
n & 因自旋引起的宽泛（“模糊”）吸收$^{[37]}$\\
\hline
nn & 非常宽泛的吸收特征$^{[24]}$\\
\hline
neb & 星云的光谱混入其中$^{[37]}$\\
\hline
p & 未指定的特殊性，奇特恒星.\\
\hline
pq & 奇特光谱，类似于新星的光谱\\
\hline
q & P型仙后座轮廓\\
\hline
s & 窄（“锐利”）吸收线$^{[37]}$\\
\hline
ss & 极窄线\\
\hline
sh & 壳层星特征$^{[37]}$\\
\hline
变量 & 可变光谱特征$^{[37]}$（有时简称为“v”）\\
\hline
波长 & 弱线$^{[37]}$（也称“w”和“wk”）\\
\hline
元素符号 & 指定元素的异常强烈的光谱线$^{[37]}$\\
\hline
z & 表明在468.6 纳米处存在异常强烈的电离氦线$^{[35]}$\\
\hline
\end{tabular}
\end{table}
例如，59 Cygni 被列为光谱型 B1.5Vnne$^{[43]}$，这表示其光谱的总体分类为 B1.5V，同时还具有非常宽的吸收线以及某些发射线。
\subsection{历史}
哈佛光谱分类中字母顺序之所以显得有些“反常”，其原因在于历史演变：该体系起源于更早的塞基（Secchi）分类，并随着人们对恒星物理认识的不断加深而逐步修订和重组。
\subsubsection{塞基分类}
在 19 世纪 60—70 年代，恒星光谱学的先驱安杰洛·塞基（Angelo Secchi） 创建了塞基光谱分类体系，用于对观测到的恒星光谱进行分类。到 1866 年，他已经提出了三类恒星光谱，如下表所示$^{[44][45][46]}$。

在 19 世纪末（1890 年代后期），这一分类体系逐渐被哈佛光谱分类所取代；后者正是本文其余部分所讨论的现代恒星光谱分类基础$^{[47][48][49]}$。
\begin{table}[ht]
\centering
\caption{请输入表格标题}\label{tab_Stcl2}
\begin{tabular}{|c|c|}
\hline
\textbf{等级编号} & \textbf{塞奇等级描述}\\
\hline
塞奇等级I & 白蓝相间的恒星，具有宽阔而强烈的氢线，例如织女星和牛郎星。这包括现代A型和早期F型。\\
\hline
塞奇等级I（猎户座亚型）& $<b$一种属于Secchi I类的亚型，其光谱中以窄线取代了宽带，例如参宿四和大角星。按现代术语来说，这对应于早期B型恒星\\
\hline
塞奇等级II & 黄色恒星——氢线较弱，但金属谱线明显，例如太阳, 大角星和五车二。这包括现代的G型和K型恒星，以及晚期F型恒星。\\
\hline
塞奇等级III & Orange to red stars with complex band spectra, such as Betelgeuse and Antares.
这对应于现代的M类。\\
\hline
塞奇四等 & In 1868, he discovered carbon stars, which he put into a distinct group:$^{[50]}$带有显著碳谱带和谱线的红色恒星，对应于现代的C类和S类。\\
\hline
塞奇五等 & 1877年，他增设了第五等：$^{[51]}$发射线星，例如仙后座γ星和谢利阿克星, 属于现代的Be型。1891年，爱德华·查尔斯·皮克林提出，V类应对应于现代的O型（当时还包括沃尔夫-拉叶星）。以及行星状星云中的恒星。\\
\hline
\end{tabular}
\end{table}
用于塞基（Secchi）光谱分类的罗马数字，不应与耶基斯（Yerkes）光度分类以及拟议的中子星分类中所使用的、彼此毫不相关的罗马数字混淆。
\subsubsection{德雷珀体系}
\begin{figure}[ht]
\centering
\includegraphics[width=14.25cm]{./figures/a92bc94a88248741.png}
\caption{} \label{fig_Stcl_3}
\end{figure}
在丈夫去世后，玛丽·安娜·德雷珀（Mary Anna Draper）开始资助哈佛底片库（Harvard Plate Stacks）的建立，以及在哈佛学院天文台对这些天文底片的研究。天文台台长爱德华·C·皮克林（Edward C. Pickering）随后雇用了多位开创性的女性天文学家，她们后来被统称为“哈佛计算员”（Harvard Computers）。尽管最初设想她们会研究多个不同的天文学课题，但这项工作的一个早期重要成果，是《亨利·德雷珀恒星光谱纪念目录》（The Henry Draper Memorial Catalogue of Stellar Spectra）的第一版，该目录于 1890 年首次出版。

威廉敏娜·弗莱明（Williamina Fleming）完成了该目录第一版中大部分光谱的分类工作，她被认为共分类了 一万余颗恒星，并发现了 10 颗新星以及 200 多颗变星。在哈佛计算员，尤其是威廉敏娜·弗莱明的帮助下，亨利·德雷珀目录的最初版本被设计出来，用以取代由安杰洛·塞基（Angelo Secchi）建立的罗马数字光谱分类体系。

该目录采用了一种新的分类方案：将此前使用的塞基分类（I 至 V 类）进一步细分为更具体的类别，并用 A 到 P 的字母表示；此外，还使用字母 Q 来标记那些无法归入其他任何类别的恒星。$^{[53][54]}$弗莱明与皮克林合作，根据氢谱线强度的不同区分出了 17 个光谱类别。氢谱线强度的变化会导致恒星辐射波长的差异，从而造成颜色外观的变化。在这一体系中，A 类光谱表现出最强的氢吸收线，而 O 类光谱几乎不显示可见谱线。随着字母顺序向后推进，氢吸收线的强度逐渐减弱。后来，安妮·跳·坎农（Annie Jump Cannon）和安东尼娅·莫里（Antonia Maury）对该体系进行了修订，最终形成了哈佛光谱分类体系。$^{[56][58]}$
\subsubsection{旧哈佛体系（1897）}
1897 年，哈佛的另一位天文学家安东尼娅·莫里将塞基 I 类中的“猎户型”亚类置于塞基 I 类其余部分之前，从而在分类顺序上使现代的 B 型位于现代的 A 型之前。她是第一位这样做的人，尽管她并未使用字母标记的光谱类型，而是采用了一个由 I–XXII 共 22 个编号组成的分类序列。$^{[59][60]}$
\begin{table}[ht]
\centering
\caption{1897年哈佛系统的概要$^{[61]}$}\label{tab_Stcl3}
\begin{tabular}{|c|c|}
\hline
\textbf{分组} & \textbf{摘要}\\
\hline
I−V & 包括了从I组到V组氢吸收线强度逐渐增强的‘猎户座型’恒星\\
\hline
VI & 充当了‘猎户座型’与Secchi I型之间的中间类型\\
\hline
VII−XI & 是塞基1型恒星，其氢吸收线强度从第VII组至第XI组逐渐减弱\\
\hline
XIII−XVI & 包括塞基2型恒星，其氢吸收线强度逐渐减弱，而太阳型金属线则逐渐增强\\
\hline
XVII−XX & $<b$包括光谱线逐渐增多的塞奇3型恒星\\
\hline
XXI & 包括第4型天狼星\\
\hline
二十二 & 包括沃尔夫-拉叶星\\
\hline
\end{tabular}
\end{table}
由于最初的 22 个罗马数字分组仍不足以涵盖恒星光谱中出现的全部变化，人们又引入了三种附加划分，以进一步细化差异：在原有编号后添加小写字母，用来区分谱线外观的相对特征；这些谱线被定义为：
\begin{itemize}
\item $a$：平均宽度
\item $b$：模糊
\item $c$：锐利
\end{itemize}
安东尼娅·莫里于 1897 年发表了她自己的恒星分类目录，题为
《用 11 英寸德雷珀望远镜拍摄的亮星光谱：亨利·德雷珀纪念项目的一部分》（Spectra of Bright Stars Photographed with the 11 inch Draper Telescope as Part of the Henry Draper Memorial）。
该目录收录了 4,800 张照片以及莫里本人对 681 颗北天亮星的分析。这是首次由女性署名的天文台正式出版物**。
\subsubsection{现行哈佛体系（1912）}
1901 年，安妮·跳·坎农重新采用了字母分类法，但她舍弃了除 O、B、A、F、G、K、M、N 之外的所有字母，并按这一顺序使用；此外，她还保留了 P（用于行星状星云）和 Q（用于某些异常光谱）。她还引入了诸如 B5A（介于 B 型与 A 型之间的恒星）、F2G（从 F 型向 G 型过渡五分之一）的混合标记方式，依此类推。

最终，到 1912 年，坎农将 B、A、B5A、F2G 等类型统一改写为 B0、A0、B5、F2 等形式。这一体系基本上就是现代哈佛光谱分类系统。该系统建立在对摄影底片上恒星光谱的系统分析之上，这种方法能够将恒星发出的光转换为可读的光谱。
\subsubsection{威尔逊山分类}
一种称为威尔逊山体系的光度分类法曾被用于区分不同光度的恒星。即便在今天，这种标记方式仍偶尔可在现代光谱中见到。其主要符号包括：
\begin{itemize}
\item sd：次矮星（subdwarf）
\item d：矮星（dwarf）
\item sg：次巨星（subgiant）
\item g：巨星（giant）
\item c：超巨星（supergiant）
\end{itemize}
\subsection{光谱型}
“光谱型”（Spectral type）一词在此指恒星光谱分类；若指小行星的光谱分类，则另见“小行星光谱类型”。恒星分类体系本质上是一种分类学（taxonomy）体系，以模式恒星为基础，类似于生物学中对物种的分类：每一个类别及其子类别都由一颗或多颗标准恒星来定义，并配有相应的判别特征描述。
\subsubsection{“早型（early）”与“晚型（late）”的命名法}
恒星常被称为早型或晚型。其中，“早型”是更热的同义词，而“晚型”是更冷的同义词。

根据具体语境，“早”和“晚”既可以是绝对概念，也可以是相对概念。作为绝对概念时，“早型”通常指 O 型或 B 型恒星，有时也包括 A 型恒星；作为相对概念时，“早型”则指在同一光谱型内部温度较高的恒星，例如“早期 K 型”可能指 K0、K1、K2 或 K3。

“晚型”的用法与此类似。不加限定地使用“晚型”时，通常指 K 型和 M 型恒星；但在相对意义上，它也可用于表示在同一光谱型中相对较冷的恒星，例如“晚期 G 型”指 G7、G8 和 G9。

在这种相对用法中，“早型”对应于光谱型字母后较小的阿拉伯数字，“晚型”则对应于较大的数字。

这种看似晦涩的术语源自 19 世纪末的一种恒星演化模型残余。该模型认为恒星的能量来自引力收缩，即开尔文–亥姆霍兹机制；而如今已知，这一机制并不适用于主序星。按照该假设，恒星会在生命初期是非常炽热的“早型星”，随后逐渐冷却，演化为“晚型星”。然而，这一机制给出的太阳年龄远小于地质记录所显示的真实年龄，并在发现恒星能量来源于核聚变之后被彻底否定。尽管如此，“早型”和“晚型”这两个术语仍被沿用下来，尽管它们所依托的理论早已被淘汰。
\subsubsection{O 型（Class O）}
\begin{figure}[ht]
\centering
\includegraphics[width=8cm]{./figures/2338d2de317ddeae.png}
\caption{一颗假想的 O5V 型恒星的光谱*} \label{fig_Stcl_4}
\end{figure}
O 型恒星极其炽热且光度极高，其辐射能量的大部分位于紫外波段。它们是所有主序星中最为稀少的一类。在太阳邻域的主序星中，大约每 3,000,000 颗才有 1 颗（约 0.00003\%）属于 O 型星 $^{[c][11]}$。一些质量最大的恒星就位于这一光谱型内。O 型恒星周围的环境往往十分复杂，这也使得其光谱测量变得困难。

早期对 O 型光谱的定义基于 He II $\lambda4541$与 He I $\lambda4471$ 谱线强度之比，其中 $\lambda$ 表示辐射波长。光谱型 O7 被定义为这两条谱线强度相等的位置，而在更早的亚型中，He I 谱线会逐渐减弱。按最初定义，O3 型对应于该谱线完全消失的情形，不过在现代技术条件下仍可极其微弱地观测到它。因此，现代定义改为采用氮谱线 N IV $\lambda4058$ 与 N III $\lambda\lambda4634$–$40$–$42$的强度比 $^{[74]}$。

O 型恒星的光谱以 He II 吸收线（有时也伴随发射线）为主，同时具有显著的电离元素谱线（如 Si IV、O III、N III、C III）以及中性氦谱线，这些氦谱线在 O5 到 O9 之间逐渐增强；还可见明显的氢巴耳末系谱线，但其强度不及更晚型恒星。质量更高的 O 型恒星由于其恒星风速度极高（可达 2,000 km/s），难以保留厚重的大气层。正因为质量巨大，O 型恒星的核心温度极高，氢燃料消耗得非常迅速，因此它们往往是最早离开主序带的恒星。

当 MKK 分类方案于 1943 年首次提出时，O 型仅包含 O5–O9.5 这些亚型 $^{[75]}$。随后，该方案在 1971 年扩展到 O9.7 $^{[76]}$，在 1978 年扩展到 O4$^{[77]}$，并在之后又引入了 O2、O3 以及 O3.5 等新的亚型 $^{[78]}$。

光谱标准星示例 $^{[72]}$：
\begin{itemize}
\item O7V —— 麒麟座 S（S Monocerotis）
\item O9V —— 蝎虎座 10（10 Lacertae）
\end{itemize}
\subsubsection{B 型恒星（Class B）}
\begin{figure}[ht]
\centering
\includegraphics[width=8cm]{./figures/a4b2537f0543257b.png}
\caption{假想的 B3V 型恒星的光谱} \label{fig_Stcl_5}
\end{figure}
B 型恒星非常明亮，呈蓝色。它们的光谱中包含中性氦谱线，在 B2 亚型中最为显著，同时还具有中等强度的氢谱线。由于 O 型和 B 型恒星能量极高，它们的寿命相对较短。因此，在其一生中发生显著动力学相互作用的概率很低，除了一些“逃逸星”（runaway stars）之外，它们通常不会远离自身形成的区域。

最初，O 型向 B 型的过渡被定义为 He II λ4541 谱线消失的界限。然而，借助现代观测设备，这条谱线在早期 B 型恒星中仍然可以观测到。如今，对于主序星而言，B 型恒星的定义改为依据 He I 紫端谱线的强度，其最大强度对应于 B2 型。对于超巨星，则使用硅谱线进行判定：Si IV λ4089 与 Si III λ4552 的谱线表明其为早期 B 型；在中期 B 型中，后者相对于 Si II λλ4128–30 的强度是关键判据；而在晚期 B 型中，则以 Mg II λ4481 相对于 He I λ4471 的强度作为区分标准。${}^{[74]}$

这些恒星往往分布在其起源的 OB 星协中，而 OB 星协通常与巨分子云相关联。猎户座 OB1 星协占据了银河系一条旋臂的相当大一部分，并包含猎户座中许多较亮的恒星。在太阳邻域的主序星中，大约每 800 颗恒星中就有 1 颗（约 0.125\%）是 B 型主序星。${}^{[c][11]}$B 型恒星相对少见，其中距离我们最近的是轩辕十四（Regulus），约 80 光年之外。${}^{[79]}$

一些质量很大但并非超巨星的恒星被称为 Be 星。这类恒星的光谱中可观测到一条或多条巴耳末系发射线，其中由恒星向外投射的与氢相关的电磁辐射系列尤为引人关注。一般认为，Be 星具有异常强烈的恒星风、较高的表面温度，并且由于其以异常快的速度自转，会发生显著的恒星质量流失。${}^{[80]}$

另一类被称为 B[e] 星（或因排版原因写作 B(e) 星）的天体，具有显著的中性或低电离态发射线。这些谱线被认为涉及禁戒跃迁机制，即发生了在现有量子力学理解下通常不被允许的过程。

典型光谱标准星：${}^{[72]}$
\begin{itemize}
\item B0V —— 猎户座 υ 星（Upsilon Orionis）
\item B0Ia —— 参宿二（Alnilam）
\item B2Ia —— 猎户座 χ² 星（Chi² Orionis）
\item B2Ib —— 仙王座 9 星（9 Cephei）
\item B3V —— 北斗七星之七·摇光（Alkaid）
\item B3V —— Haedus
\item B3Ia —— 大犬座 ο² 星（Omicron² Canis Majoris）
\item B5Ia —— 船尾座 ε 星（Aludra）
\item B8Ia —— 参宿七（Rigel）
\end{itemize}
\subsubsection{A 型恒星（Class A）}
\begin{figure}[ht]
\centering
\includegraphics[width=8cm]{./figures/a2e592b0bb3c9ac1.png}
\caption{一颗假想的 A5V 型恒星的光谱} \label{fig_Stcl_6}
\end{figure}
A 型恒星是在肉眼可见恒星中较为常见的一类，颜色呈白色或蓝白色。它们具有很强的氢谱线，在 A0 型达到最强；同时还具有电离金属谱线（如 Fe II、Mg II、Si II），其强度在 A5 型达到最大。到这一阶段，Ca II 谱线的增强也尤为明显。在太阳邻域的主序星中，大约每 160 颗恒星中就有 1 颗（约 0.625\%）是 A 型恒星，${}^{[c][11]}$其中在 15 秒差距范围内就包含 9 颗 A 型恒星。$^{[81]}$

典型光谱标准星：${}^{[72]}$
\begin{itemize}
\item A0Van —— 天仓五（Phecda）
\item A0Va —— 织女星（Vega）
\item A0Ib —— 狮子座 η 星（Eta Leonis）
\item A0Ia —— HD 21389
\item A1V —— 天狼星 A（Sirius A）
\item A2Ia —— 天津四（Deneb）
\item A3Va —— 南鱼座 α 星（Fomalhaut）
\end{itemize}
\subsubsection{F 型恒星（Class F）}
\begin{figure}[ht]
\centering
\includegraphics[width=8cm]{./figures/88e11dd026c5811b.png}
\caption{一颗假想的 F5V 型恒星的光谱} \label{fig_Stcl_7}
\end{figure}
F 型恒星的光谱中，钙的 Ca II H 与 K 线逐渐增强；到晚期 F 型时，中性金属线（如 Fe I、Cr I）开始超过电离金属线。其光谱特征表现为氢线较弱、电离金属线较弱。它们的颜色呈白色。在太阳邻域的主序星中，大约 1/33（3.03\%）属于 F 型恒星$^{[c][11]}$，其中在 20 光年范围内仅有一颗，即南河三 A（Procyon A）$^{[82]}$。

\textbf{典型光谱标准星：}$^{[72][83][84][85][86]}$
\begin{itemize}
\item F0IIIa —— Adhafera
\item F0Ib —— Arneb
\item F1V —— 37 Ursae Majoris（大熊座 37）
\item F2V —— 78 Ursae Majoris（大熊座 78）
\item F7V —— Iota Piscium（双鱼座 ι）
\item F9V —— Zavijava
\item F9V —— HD 10647
\end{itemize}
\subsubsection{G 型}
\begin{figure}[ht]
\centering
\includegraphics[width=8cm]{./figures/d758de7d2c3b6f51.png}
\caption{假想的 G5V 型恒星光谱} \label{fig_Stcl_8}
\end{figure}
G 型恒星（包括太阳在内）具有显著的 Ca II 的 H、K 吸收线，在 G2 亚型时最为明显。它们的氢吸收线比 F 型恒星更弱，但除电离金属线外，还出现中性金属线。在 CN 分子的 G 带处存在一个显著的谱带峰值。G 型主序星约占太阳邻域主序星的 7.5\%，也就是大约每 13 颗主序星中就有 1 颗是 G 型恒星。在 10 pc 范围内共有 21 颗 G 型恒星。

G 型还包含所谓的“黄色演化空隙”（Yellow Evolutionary Void）。超巨星在演化过程中常在 O 或 B 型（蓝色）与 K 或 M 型（红色）之间摆动，但它们在不稳定的黄色超巨星阶段停留的时间通常很短。

\textbf{光谱标准示例：}$^{[72]}$
\begin{itemize}
\item G0V —— Chara
\item G0IV —— Muphrid
\item G0Ib —— Sadalsuud
\item G2V —— 太阳
\item G5V —— Kappa¹ Ceti
\item G5IV —— Mu Herculis
\item G5Ib —— 9 Pegasi
\item G8V —— 61 Ursae Majoris
\item G8IV —— Alshain
\item G8IIIa —— Kappa Geminorum
\item G8IIIab —— Vindemiatrix
\item G8Ib —— Mebsuta
\end{itemize}
\subsubsection{K 型}
\begin{figure}[ht]
\centering
\includegraphics[width=8cm]{./figures/2bdd33a2f5bc7e0e.png}
\caption{假想的 K5V 型恒星光谱} \label{fig_Stcl_9}
\end{figure}
K 型恒星是呈橙色的恒星，温度略低于太阳。它们约占太阳邻域主序星的 12\%$^{[c][11]}$。此外还存在 K 型巨星，其范围从像 RW Cephei 这样的超巨星，到巨星和超巨星（如大角星 Arcturus）；而像半人马座 $\alpha B$这样的橙矮星则属于主序星。

K 型恒星的氢吸收线极其微弱，甚至可能完全不存在，其光谱主要由中性金属线（如 Mn I、Fe I、Si I）构成。在晚 K 型恒星中，会开始出现氧化钛（TiO）分子带。主流理论（基于较低的有害辐射水平以及恒星寿命较长）因此认为：若行星上的生命形式与地球生命直接类比，这类恒星由于其宜居带较宽、且相较于拥有最宽宜居带的恒星其高能有害辐射阶段显著更少，可能是孕育高度演化生命的最优候选恒星之一$^{[88][89]}$。

\textbf{光谱标准示例：}$^{[72]}$
\begin{itemize}
\item K0V —— Alsafi
\item K0III —— Pollux
\item K0III —— Aljanah
\item K2V —— Ran
\item K2III —— 蛇夫座 κ（Kappa Ophiuchi）
\item K3III —— 牧夫座 ρ（Rho Boötis）
\item K5V —— 天鹅座 61 A（61 Cygni A）
\item K5III —— Eltanin
\end{itemize}
\subsubsection{M 型}
\begin{figure}[ht]
\centering
\includegraphics[width=8cm]{./figures/cdd928e0cdd68726.png}
\caption{假想的 M5V 型恒星光谱} \label{fig_Stcl_10}
\end{figure}
M 型恒星（Class M）是目前为止数量最多的一类恒星。在太阳邻域的主序星中，大约 76\%属于 M 型恒星 $^{[c][f][11]}$。然而，M 型主序星（即红矮星）的光度极低，在一般情况下没有一颗亮到可以用肉眼直接观测，除非在极其特殊的条件下。已知最亮的 M 型主序星是拉卡伊 8760（Lacaille 8760），光谱型为 M0V，视星等为 6.7（而在良好观测条件下，通常认为肉眼可见的极限星等约为 6.5），因此几乎不可能再发现更亮的同类恒星。

尽管大多数 M 型恒星是红矮星，但银河系中一些体积最大的已知超巨星同样属于 M 型，例如大犬座 VY（VY Canis Majoris）、仙王座 VV（VV Cephei）、心宿二（Antares）和参宿四（Betelgeuse）。此外，一些体积较大、温度较高的棕矮星也属于晚期 M 型，通常位于 M6.5 到 M9.5 的范围内。

M 型恒星的光谱包含氧化物分子的谱线（在可见光波段中尤以 TiO 为主）以及各种中性金属的谱线，但通常不出现氢的吸收线。在 M 型恒星中，TiO 分子带往往非常强烈，到了大约 M5 型时，几乎主导了整个可见光光谱。到更晚的 M 型阶段，还会出现二价钒氧化物（VO）的分子带。

\textbf{典型光谱标准星：}$^{[72]}$
\begin{itemize}
\item M3V —— 格利泽 581（Gliese 581）
\item M0IIIa —— 米拉赫（Mirach）
\item M2III —— 飞马座 χ（Chi Pegasi）
\item M1–M2Ia–Iab —— 参宿四（Betelgeuse）
\item M2Ia —— 仙王座 μ（Mu Cephei，“赫歇尔的石榴星”）
\end{itemize}
\subsection{扩展光谱类型}
随着新类型恒星的不断发现，天文学中已经引入了若干新的光谱类型 $^{[90]}$。
\subsubsection{高温蓝色发射星类别}
\begin{figure}[ht]
\centering
\includegraphics[width=8cm]{./figures/d8662d48503162d5.png}
\caption{UGC 5797：一座发射线星系，其中正在形成质量巨大、明亮的蓝色恒星$^{[91]}$} \label{fig_Stcl_11}
\end{figure}
一些极高温、呈蓝色的恒星的光谱中，会出现显著的发射线，这些发射线主要来自碳或氮，有时也来自氧。

\textbf{WR 类（或 W 类）：沃尔夫–拉叶星（Wolf–Rayet）}
\begin{figure}[ht]
\centering
\includegraphics[width=6cm]{./figures/4033fec8a3a32df2.png}
\caption{哈勃空间望远镜拍摄的M1-67 星云及其中心的沃尔夫–瑞叶星 WR 124的图像} \label{fig_Stcl_12}
\end{figure}
曾一度被归入 O 型恒星的W 型（或 WR 型）沃尔夫–瑞叶星$^{[92]}$，其显著特征是光谱中缺乏氢谱线。相反，它们的光谱主要由高度电离的氦、氮、碳，有时还有氧的宽发射线所主导。人们认为，这类恒星大多是演化晚期的超巨星，其外层氢包层已被强烈的恒星风剥离，从而直接暴露出炽热的氦壳层。
WR 型又可根据光谱中氮与碳发射线的相对强度（也反映其外层成分）进一步细分$^{[42]}$。

WR 光谱的分类如下$^{[93][94]}$：

\begin{itemize}
\item WN$^{[42]}$ —— 光谱由 N III–V 和 He I–II 谱线主导

\item WNE（WN2–WN5，部分 WN6）：较热，称为“早型”
\item WNL（WN7–WN9，部分 WN6）：较冷，称为“晚型”
\item 有时使用扩展的 WN10 与 WN11 类型，用于 Ofpe/WN9 恒星$^{[42]}$
\item h标记（如 WN9h）：表示存在氢发射线 ha 标记（如 
WN6ha）：表示同时存在氢的发射与吸收
\item WN/C —— 兼具 WN 特征并出现强 C IV 谱线，介于 WN 与 WC 之间$^{[42]}$
\item WC$^{[42]}$ —— 光谱中具有强烈的 C II–IV 谱线
\item WCE（WC4–WC6）：较热，“早型”
\item WCL（WC7–WC9）：较冷，“晚型”
\item WO（WO1–WO4）—— 具有强 O VI 谱线，极为罕见，可视为 WCE 类向极端高温（可达 200 kK 甚至更高）的延伸
\end{itemize}
尽管大多数行星状星云的中心星（CSPNe）呈现 O 型光谱$^{[95]}$，但约 10\%的此类中心星缺乏氢，并表现出 WR 型光谱$^{[96]}$。这些恒星属于低质量恒星，为与高质量的沃尔夫–瑞叶星区分，其光谱类型用方括号标注，例如 [WC]。其中大多数呈现[WC]光谱，少数为[WO]，而[WN]则极为罕见。

\textbf{斜线星}

斜线星是光谱中同时呈现 WN 型特征谱线的 O 型恒星。名称“斜线（slash）”源于其印刷的光谱类型中包含一个斜杠（例如 Of/WNL）$^{[74]}$。

在这类光谱中还发现了一个次级族群：温度更低、处于“中间态”的恒星，记作 Ofpe/WN9$^{[74]}$。这些恒星也曾被称为 WN10 或 WN11，但随着人们认识到它们与其他沃尔夫–瑞叶星在演化路径上的差异，这种称谓已不再常用。近年来，对更为罕见恒星的发现进一步将斜线星的范围扩展至 O2–3.5If/WN5–7，其温度甚至高于最初定义的斜线星 $^{[97]}$。

\textbf{磁性 O 型恒星}

这是一类具有强磁场的 O 型恒星，其指定标记为 Of?p $^{[74]}$。
\subsubsection{冷红矮星与褐矮星类别}
新的光谱类型 L、T 和 Y 被引入，用于对冷恒星的红外光谱进行分类。这一体系既包括红矮星，也包括在可见光波段极其暗弱的褐矮星 $^{[98]}$。

褐矮星是不发生氢核聚变的天体，它们会随着年龄增长而逐渐冷却，因此会**向更晚的光谱类型演化。褐矮星在形成初期通常呈现 M 型光谱，随后会依次冷却并经历 L、T、Y 光谱类型，而且质量越小，冷却速度越快；质量最高的褐矮星在宇宙当前年龄内甚至还未冷却到 Y 型，乃至 T 型。正因如此，对于某些质量与年龄组合的 L–T–Y 类型，其有效温度与光度之间会出现不可区分的重叠，因此无法为这些类型给出明确的温度或光度取值$^{[10]}$。
\subsubsection{L 型}
\begin{figure}[ht]
\centering
\includegraphics[width=6cm]{./figures/a812750df0497cb6.png}
\caption{L 型矮星的艺术想象图} \label{fig_Stcl_13}
\end{figure}
L 型矮星之所以被如此命名，是因为它们比 M 型恒星更冷，而字母 L 在字母顺序上是最接近 M 的尚未使用字母。这类天体中，有一部分质量足够大，能够维持氢核聚变，因此严格来说仍属于恒星；但绝大多数质量低于恒星下限，因而属于褐矮星。它们在可见光下呈现出非常深的红色，而在红外波段最为明亮。其大气温度足够低，使得金属氢化物以及碱金属的光谱线在其光谱中尤为显著 $^{[99][100][101]}$。

由于巨星具有较低的表面引力，含有 TiO 和 VO 的凝结物无法形成。因此，除矮星之外的 L 型恒星在孤立环境中不可能形成。然而，通过恒星碰撞的方式，理论上仍有可能形成 L 型超巨星；一个例子是 V838 Monocerotis，它在发生高光度红新星爆发的高峰期曾表现出类似特征。
\subsubsection{T 型}
\begin{figure}[ht]
\centering
\includegraphics[width=6cm]{./figures/b57b3066df5ac4fb.png}
\caption{T 型矮星的艺术想象图} \label{fig_Stcl_14}
\end{figure}
T 型矮星是低温褐矮星，其表面温度约在 550–1,300 K（即 277–1,027 °C；530–1,880 °F）之间。它们的辐射峰值位于红外波段，并且其光谱中甲烷（$CH{_4}$）特征非常显著$^{[99][100]}$。

对原行星盘（proplyds）——即星云中由气体团块构成、恒星和行星系统由此形成的原行星盘——数量的研究表明，银河系中的恒星数量可能比先前推测的高出数个数量级。一种理论认为，这些原行星盘之间处于一种“竞赛”状态：最先形成的原行星会演化为原恒星，而原恒星是极为剧烈的天体，会扰乱其周围的其他原行星盘，并剥离它们的气体。这些“受害”的原行星盘随后很可能演化为主序星，或成为 L 型和 T 型褐矮星——而这类天体对我们来说几乎不可见 $^{[102]}$。
\subsubsection{Y 型}
\begin{figure}[ht]
\centering
\includegraphics[width=6cm]{./figures/02ecc07011945347.png}
\caption{Y 型矮星的艺术想象图} \label{fig_Stcl_15}
\end{figure}
Y 型光谱级的褐矮星比 T 型光谱级的褐矮星更冷，其光谱特征也与后者在定性上存在差异。截至 2013 年 8 月，共有 17 个天体被归入 Y 型光谱级 $^{[103]}$。尽管这类褐矮星已经被建立模型研究 $^{[104]}$，并且借助广域红外巡天探测器（WISE）在 40 光年范围内被探测到 $^{[90][105][106][107][108]}$，但目前尚未建立一个定义明确的光谱序列，也缺乏公认的原型天体。尽管如此，已有若干天体被提议划分为 Y0、Y1 和 Y2 光谱型 $^{[109]}$。

这些候选 Y 型天体的光谱在约 1.55 微米处显示出吸收特征 $^{[110]}$。Delorme 等人提出，该特征源于氨（NH(_3)）的吸收，并认为这一特征可作为 T–Y 光谱过渡的指示性标志 $^{[110][111]}$。事实上，这一氨吸收特征已成为定义 Y 型光谱级所采用的主要判据 $^{[109]}$。然而，该特征很难与水和甲烷的吸收区分开来 $^{[110]}$，也有其他研究者指出，将某些天体直接归类为 Y0 仍为时过早 $^{[112]}$。

最新被提议归入 Y 型光谱级的褐矮星是 WISE 1828+2650，它是一颗 >Y2 型褐矮星，其最初估计的有效温度约为 300 K，与人体温度相当 $^{[105][106][113]}$。然而，后续的视差测量表明，其光度与“低于约 400 K”的温度假设并不一致。目前已知最冷的 Y 型褐矮星是 WISE 0855−0714，其温度约为 250 K，质量约为木星的 7 倍 $^{[114]}$。

Y 型褐矮星的质量范围大致为 9–25 个木星质量，但年轻天体的质量可能低于 1 个木星质量（不过它们会继续冷却并演化为行星），这意味着 Y 型天体跨越了 13 个木星质量这一氘聚变极限，而该极限正是国际天文学联合会（IAU）当前用于区分褐矮星与行星的分界线 $^{[109]}$。

\textbf{异常褐矮星}
\begin{figure}[ht]
\centering
\includegraphics[width=10cm]{./figures/03ce6de093f4f103.png}
\caption{} \label{fig_Stcl_16}
\end{figure}
年轻褐矮星由于相较于具有相同光谱型的场星而言半径更大、质量更低，因此具有较低的表面重力。这类天体通常用希腊字母进行标记：$\beta$（贝塔）表示中等表面重力，$\gamma$（伽马）表示低表面重力。

低表面重力的判据包括：**CaH、K I 和 Na I 谱线较弱，而 VO 谱线较强** $^{[117]}$。
**α（阿尔法）**表示**正常表面重力**，通常会被省略不用。
在某些情况下，**极低的表面重力**会用 **δ（德尔塔）**表示 $^{[119]}$。

后缀 **“pec”** 表示**异常（peculiar）**。这一后缀仍用于描述其他不寻常的光谱特征，是一个综合性标记，可指示低表面重力、亚矮星或未分辨的双星等多种性质 $^{[120]}$。

前缀 sd 表示亚矮星（subdwarf），且仅用于冷亚矮星。该前缀意味着低金属丰度，并指示其运动学性质更接近于银晕星而非银盘星 $^{[116]}$。亚矮星在颜色上通常比银盘天体更偏蓝 $^{[121]}$。

red（红）后缀用于描述颜色偏红、但年龄较老的天体。这种红色并不被解释为低表面重力，而被认为源于较高的尘埃含量$^{[118][119]}$。blue（蓝）后缀则描述近红外颜色偏蓝、且无法用低金属丰度解释的天体。其中一部分可解释为 L+T 型双星，而另一些并非双星，例如 2MASS J11263991−5003550，其蓝色特征被认为源于稀薄和/或大颗粒云层$^{[119]}$。
\subsubsection{晚期巨型碳星光谱级}
碳星（carbon stars）是指其光谱显示出碳元素生成的恒星，而碳正是三 α 氦聚变反应的副产物。随着碳丰度的增加，并伴随部分 s 过程重元素的生成，这些恒星的光谱会逐渐偏离通常的晚期光谱型 G、K 和 M。对于富碳恒星，其对应的等效光谱级为 S 型和 C 型。

其中的巨星通常被认为是自身产生了碳元素；但该类中也存在一些双星系统，其异常的大气成分被推测是在伴星仍为碳星、而后演化为白矮星之前，通过物质转移获得的。

\textbf{C 型恒星}
\begin{figure}[ht]
\centering
\includegraphics[width=6cm]{./figures/ef4fd3c47b2bcae4.png}
\caption{碳星 R Sculptoris 的图像，展示了其引人注目的螺旋结构} \label{fig_Stcl_17}
\end{figure}
最初被归类为 **R** 型和 **N** 型恒星，它们也被统称为**碳星**。这类恒星通常是生命晚期的红巨星，其大气中碳元素含量过量。早期的 R 型与 N 型分类体系大致与常规光谱分类中从中等 **G 型** 到晚期 **M 型** 平行存在。近年来，这一体系被重新整合为统一的碳星分类 **C 型**，其中 **N0** 大致对应于 **C6**。

另一类较冷的碳星子类是 C–J 型星，其特征是在含有(^{12}\mathrm{CN}) 分子的同时，还显著富含$^{13}\mathrm{CN}$分子$^{[122]}$。目前已知少数主序碳星，但绝大多数已确认的碳星都是巨星或超巨星。主要的若干子类包括：
\begin{itemize}
\item C-R —— 原本是独立的 R 型，代表与晚期 G 型至早期 K 型恒星对应的碳星；
\item C-N —— 原本是独立的 N 型，代表与晚期 K 型至 M 型恒星对应的碳星；
\item C-J —— 冷 C 型星的一个亚型，具有较高的$^{13}\mathrm{C}$ 含量；
\item C-H —— C-R 星在 II 星族中的对应类型；
\item C-Hd —— 缺氢碳星，类似于晚期 G 型超巨星，但其光谱中额外出现 CH与 $C{_2}$ 分子带。
\end{itemize}

\textbf{S 型}

S 型恒星在 M 型恒星与碳星之间形成一条连续谱系。与 M 型最为相似的 S 型星，其光谱中具有强烈的 ZrO 吸收带，类似于 M 型恒星中的 TiO 吸收带；而与碳星最相似的 S 型星，则表现为强烈的钠 D 线和较弱的$C{_2}$分子带 $^{[123]}$。

S 型恒星中锆及其他由 s 过程产生的元素含量异常丰富，其碳与氧的相对丰度介于 M 型恒星与碳星之间。与碳星类似，几乎所有已知的 S 型恒星都处在渐近巨星支（AGB）阶段。

S 型光谱类型用字母 S 加上 0–10 之间的数字表示，该数字对应恒星的温度，并大致遵循 M 型巨星的温标。最常见的类型为 S3–S5。非标准的 S10 仅在天鹅座 χ（Chi Cygni）处于极小光度时使用过。

基础分类后通常会附加丰度标记，常见形式包括：S2,5；S2/5；S2 Zr4 Ti2；或 S25。逗号后的数字（1–9）表示 ZrO/TiO 比值；斜杠后的数字是较新的、但不太常用的方案，用于表示 C/O 比值（1–10），其中 0 表示 MS 星；锆与钛的强度有时会被明确标注；星号后的数字（1–5）表示 ZrO 分子带的强度。

\textbf{MS 与 SC 型：碳相关的过渡类型}

位于 M 型与 S 型之间的边界恒星被称为 MS 星；类似地，位于 S 型与 C-N 型之间的边界对象被称为 SC 星（或 CS 星）。M → MS → S → SC → C-N这一序列被认为反映了渐近巨星支碳星随年龄增长、碳丰度逐渐增加的演化过程。
\subsubsection{白矮星的分类}
D 类（Degenerate，简并星）是现代用于白矮星的分类。白矮星是质量较低、已不再发生核聚变的恒星，其体积收缩到行星大小，并在此后缓慢冷却。D 类进一步细分为若干光谱类型：DA、DB、DC、DO、DQ、DX、DZ。这些字母与其他恒星分类体系中的字母无关，而是用来指示白矮星可见外层或大气的化学组成。

白矮星的各类型如下 $^{[124][125]}$：

\begin{itemize}
\item DA —— 富氢大气或外层，由强烈的巴耳末氢谱线指示。
\item DB —— 富氦大气，由中性氦（He I）谱线指示。
\item DO —— 富氦大气，由电离氦（He II）谱线指示。
\item DQ —— 富碳大气，由原子或分子碳谱线指示。
\item DZ —— 富金属大气，由金属谱线指示（为已废弃的白矮星光谱类型 DG、DK、DM的合并）。
\item DC —— 没有显示上述任何类型所需的强谱线。
\item DX —— 谱线不够清晰，无法归入上述任一类型。
\end{itemize}

类型字母之后通常会跟随一个数字，用以表示白矮星的表面温度。该数字是$50400 / T_{\mathrm{eff}}$的取整形式，其中$T_{\mathrm{eff}}$为以开尔文计的有效表面温度。最初该数字只取 1–9 之间的整数，但近年来开始使用分数值，以及小于 1 或大于 9 的数值（例如 IK Pegasi B 的分类为 DA1.5）$^{[124][126]}$。

当一颗白矮星同时呈现以上多种光谱特征时，可以使用两个或更多类型字母进行联合标注 $^{[124]}$。

\textbf{扩展的白矮星光谱类型}
\begin{figure}[ht]
\centering
\includegraphics[width=6cm]{./figures/004e17f144e99498.png}
\caption{由哈勃望远镜分辨开的天狼星 A 与天狼星 B（其中天狼星 B 为 DA2 型白矮星）。} \label{fig_Stcl_18}
\end{figure}
\begin{itemize}
\item DAB —— 富含氢和氦的白矮星，光谱中显示中性氦谱线
\item DAO —— 富含氢和氦的白矮星，光谱中显示电离氦谱线
\item DAZ —— 富含氢、同时含有金属元素的白矮星
\item DBZ —— 富含氦、同时含有金属元素的白矮星
\end{itemize}
白矮星所使用的一套光谱特异性符号与其他类型恒星所用的符号体系不同。
\begin{table}[ht]
\centering
\caption{请输入表格标题}\label{tab_Stcl4}
\begin{tabular}{|c|c|}
\hline
\textbf{代码} & \textbf{恒星的光谱特征}\\
\hline
\textbf{P} & 具有可探测偏振的磁性白矮星\\
\hline
\textbf{E} & 存在发射线\\
\hline
\textbf{H} & 无可探测偏振的磁性白矮星\\
\hline
\textbf{V} & 变星\\
\hline
\textbf{PEC} & 存在光谱异常\\
\hline
\end{tabular}
\end{table}
\subsubsection{发光蓝变星}
发光蓝变星（LBV）是一类罕见、质量极大且处于演化晚期的恒星，其光谱和亮度会表现出不可预测、且有时非常剧烈的变化。在其“静止态”下，它们通常与B 型恒星相似，但会出现一些异常的光谱线；而在爆发期，它们更接近 F 型恒星，其表面温度显著降低。许多文献将 LBV 视为一种独立的光谱类型。$^{[127][128]}$
\subsubsection{非单一恒星天体的光谱类型：P 类与 Q 类}
最后，P 类和 Q 类是安妮·坎农（Annie Jump Cannon）在为《亨利·德雷珀星表》建立分类体系时遗留下来的类别。这两类偶尔用于某些并非对应单一恒星的天体：P 型天体：位于行星状星云中的恒星（通常是年轻的白矮星，或贫氢的 M 型巨星）；Q 型天体：新星。$^{[citation needed]}$
\subsection{恒星遗迹}
恒星遗迹是与恒星死亡相关的天体。这一类别包括白矮星；正如 D 类所采用的截然不同的分类方案所显示的那样，恒星遗迹很难被纳入 MK 分类体系之中。

MK 体系所依托的赫罗图（赫茨普朗—罗素图）本质上是观测型的，因此这些遗迹难以在图中标注，甚至无法放置。年老的中子星体积相对较小、温度较低，若能标注，通常会落在图的最右侧。行星状星云具有高度的动力学特征，随着其前身恒星过渡到白矮星分支，其亮度会迅速衰减；若绘制出来，行星状星云应位于赫罗图右上象限的右侧。黑洞本身不发射可见光，因此不会出现在赫罗图上。$^{[129]}$

曾有人提出一种使用罗马数字的中子星分类体系：I 型：质量较小、冷却速率较低的中子星；II 型：质量更大、冷却速率更高的中子星；III 型（拟议）：质量更大、冷却更快的中子星（可能是奇异星候选）。$^{[130]}$中子星质量越大，其中微子通量越高；中微子会带走大量热能，使得孤立中子星在仅仅几年内就能从数十亿开尔文降至约一百万开尔文。需要注意的是，这一拟议的中子星分类不应与早期的塞基光谱分类或耶基斯光度等级相混淆。
\subsection{被替代的光谱类型}
若干在 20 世纪中期用于非标准恒星的光谱类型，已在恒星分类体系的修订过程中被取代；它们仍可能出现在旧版星表中：R 类与 N 类已被并入新的 C 类，分别成为 C-R 与 C-N。
\subsection{恒星分类、宜居性与生命搜索}
尽管人类未来或许能够在任何类型的恒星环境中定居，但本节关注的是生命在其他恒星周围出现的概率。

稳定性、光度与寿命都是恒星宜居性的关键因素。人类目前仅知一个承载生命的恒星——G 型的太阳，它具有较高的重元素丰度且亮度变化小。此外，太阳系也不同于许多恒星系统，因为它只包含一颗恒星（见双星系统的宜居性）。

在这些约束条件以及样本量仅为 1 的经验性限制下，被预测可能支持生命的恒星范围受制于少数因素。就主序星而言，质量超过太阳 1.5 倍的恒星（O、B、A 型）演化过快，不利于高级生命的发展（以地球为参照）。而在另一端，质量低于太阳一半的矮星（M 型）很可能使其宜居带内的行星发生潮汐锁定，并伴随其他问题（见红矮星系统的宜居性）。$^{[131]}$ 尽管红矮星面临诸多不利因素，许多天文学家仍持续对其建模研究，原因在于其数量极多且寿命很长。

基于这些原因，NASA 的开普勒任务将重点搜索附近主序星周围的宜居行星，这些恒星的质量低于 A 型、但高于 M 型——因此，F、G、K 型矮星被认为是最有可能承载生命的恒星类型。$^{[131]}$
\subsection{另见}
\begin{itemize}
\item 天体照相望远镜（Astrograph）——一种望远镜类型
\item 客星（Guest star）——中国古代对灾变变星的称谓
\item 光谱指纹（Spectral signature）——材料反射或辐射随波长变化的特征
\item 恒星计数（Star count）——对恒星进行清点与统计的巡天调查
\item 恒星动力学（Stellar dynamics）——天体物理学的一个分支
\end{itemize}
\subsection{注释}\\
a.这是在以织女星（Vega）（通常被视为偏蓝的恒星）作为“白色”标准时，恒星的相对颜色。\\
b.色度在同一光谱类型内可能有显著差异；例如，太阳（G2 型）呈白色，而 G9 型恒星则呈黄色。\\
c.这些比例指的是绝对星等亮于 16 等的恒星所占分数；降低这一阈值会使早型星显得更加稀少，而增加的数量通常主要落在 M 型。这些比例在计算时忽略了总计列中的 800，因为实际数目加总为 824。\\
d.从严格意义上说，白矮星已不再是“存活”的恒星，而是已熄灭恒星的“遗骸”；因此，其分类采用了一套不同于进行元素燃烧的“活恒星”的光谱类型体系。\\
e.当用于 A 型恒星时，该术语指的是异常强的金属光谱线。\\
f.若将所有恒星都计入，这一比例会上升到 78.6\%（见上文注释）。
\subsection{参考文献}
\begin{enumerate}
\item “摩根-基南光度等级 | 宇宙”.astronomy.swin.edu.au。检索日期：2022年8月31日
\item 奥康奈尔（2023年3月27日）。“星等与颜色系统”（PDF）。加州理工学院天文学511课程于2023年3月28日存档自原始版本。检索日期：2023年3月27日
\item 杰弗里，C. S.；维尔纳，K.；基尔肯尼，D.；米沙尔斯基，B.；莫纳根，I.；斯诺登，E. J.（2023）。“利用SALT发现的炽热白矮星和前白矮星”. 皇家天文学会月报. 519 (2): 2321–2330. arXiv:2301.03550. doi:10.1093/mnras/stac3531.
\item 哈贝茨，G. M. H. J.；海因泽，J. R. W.（1981年11月）。“主序星的实测辐射度校正”。天文学与天体物理学增刊系列. 46：193–237（表VII和VIII）。文献标识码:1981A&AS...46..193H.——光度由Mbol值推算得出，采用Mbol(☉)=4.75.
\item 魏德纳，卡斯滕；文克，约里克·S.（2010年12月）。“恒星的质量与O型星的质量差异”。天文学与天体物理学. 524. A98。arXiv:1010.2204. Bibcode:2010A&A...524A..98W. doi:10.1051/0004-6361/201014491. S2CID 118836634.
\item 慈善，米切尔。“星星是什么颜色的？”.Vendian.org。检索日期：2006年5月13日
\item “恒星的颜色”。澳大利亚望远镜国家设施。2018年10月17日。
\item 摩尔，帕特里克（1992）。《吉尼斯天文学大全：事实与壮举》（第4版）。吉尼斯。ISBN 978-0-85112-940-2.
\item “恒星的颜色”。澳大利亚望远镜推广与教育项目。2004年12月21日。于2013年12月3日从原文存档。检索日期：2007年9月26日。— 解释了颜色感知差异的原因。
\item 巴拉夫，I.；沙布里耶，G.；巴尔曼，T. S.；阿拉尔，F.；豪斯希尔德特，P. H.（2003年5月）。“冷棕矮星与太阳系外巨行星的演化模型：以HD 209458为例”。天文学与天体物理学. 402 (2): 701–712. arXiv:astro-ph/0302293. 文献标识码:2003A&A...402..701B. DOI:10.1051/0004-6361:20030252. S2CID 15838318.
\item 莱德鲁，格伦（2001年2月）。“真正的星空”。加拿大皇家天文学会期刊. 95: 32. 文献标识码:2001JRASC..95...32L.
\item “恒星的光谱分类（OBAFGKM）”.www.eudesign.com。检索日期：2019年4月6日
\item “哈佛光谱分类法的记忆口诀”。检索日期：2025年6月10日.
\item “AST 101：伟大的记忆术竞赛”。检索日期：2025年6月10日。
\item 索塔，A.；迈斯·阿佩亚尼兹，J.；莫雷尔，N. I.；巴尔瓦，R. H.；沃尔本，N. R.；等（2014年3月）。“银河系O型恒星光谱巡天（GOSSS）。II. 南方亮星”。天体物理学报增刊系列. 211 (1). 10. arXiv:1312.6222. 文献编码:2014ApJS..211...10S. DOI:10.1088/0067-0049/211/1/10. S2CID 118847528.
\item 菲利普斯，肯尼思·J·H.（1995）。太阳指南。剑桥大学出版社。第47–53页。ISBN978-0-521-39788-9
\item 拉塞尔，亨利·诺里斯（1914年3月）。“恒星光谱与其他特征之间的关系”。大众天文学。第22卷。第275–294. 文献标识码:1914PA.....22..275R.
\item 萨哈，M. N.（1921年5月）。“关于恒星光谱的物理理论”. 伦敦皇家学会会刊·A辑. 99 (697): 135–153. 文献标识码:1921RSPSA..99..135S. DOI:10.1098/rspa.1921.0029.
\item 佩恩，塞西莉亚·海伦娜（1925）。恒星大气层：对恒星反转层高温观测研究的贡献（博士）。拉德克利夫学院。文献标识码:1925PhDT.........1P.
\item 宇宙、物理与（2013年6月14日）。“耶克斯光谱分类”。物理与宇宙。检索日期：2022年8月31日
\item 伦敦大学学院（2018年11月30日）。“MKK与修订版MK地图集”。伦敦大学学院天文台（UCLO）。检索日期：2022年8月31日
\item 摩根，威廉·威尔逊；基南，菲利普·蔡尔兹；凯尔曼，埃迪丝（1943）。恒星光谱图集，附光谱分类概要。芝加哥大学出版社。天文学文献标识符:1943assw.book.....M. OCLC 1806249.
\item 摩根，威廉·威尔逊；基南，菲利普·蔡尔斯（1973）。“光谱分类”。天文学与天体物理学年评. 11: 29–50. 文献标识码:1973ARA&A..11...29M. doi:10.1146/annurev.aa.11.090173.000333.
\item “关于光谱图集与光谱分类的说明”。斯特拉斯堡天文数据中心。检索日期：2015年1月2日
\item 卡瓦列罗-尼维斯，S. M.；尼兰，E. P.；吉斯，D. R.；华莱士，D. J.；德乔伊-伊斯特伍德，K.；等（2014年2月）。“天鹅座OB2中大质量恒星的高角分辨率巡天：哈勃空间望远镜精细导引传感器的观测结果”。天文学杂志. 147 (2). 40. arXiv:1311.5087. Bibcode:2014AJ....147...40C. doi:10.1088/0004-6256/147/2/40. S2CID 22036552.
\item 普林贾，R. K.；马萨，D. L.（2010年10月）。“B型超巨星风中广泛团聚现象的特征”。天文学与天体物理学. 521. L55.arXiv:1007.2744. Bibcode:2010A&A...521L..55P. doi:10.1051/0004-6361/201015252. S2CID 59151633.
\item 格雷，大卫·F.（2010年11月）。“超巨星天鹅座γ的光球变化”. 《天文学杂志》. 140 (5): 1329–1336. 文献标识码:2010AJ....140.1329G. DOI:10.1088/0004-6256/140/5/1329.
\item 纳泽，Y.（2009年11月）。“XMM-牛顿观测的炽热恒星。I. 目录与OB型恒星的性质”。天文学与天体物理学.506(2)：1055–1064。arXiv：0908.1461Bibcode2009A&A...506.1055Ndoi10.1051/0004-6361/200912659S2CID17317459
\item 柳比姆科夫，列昂尼德·S.；兰伯特，大卫·L.；罗斯托普钦，谢尔盖·I.；拉奇科夫斯卡娅，塔玛拉·M.；波克拉德，德米特里·B.（2010年2月）。“太阳邻域中A型、F型和G型超巨星的精确基本参数”. 皇家天文学会月报. 402 (2): 1369–1379. arXiv:0911.1335. 文献标识码:2010MNRAS.402.1369L. DOI:10.1111/j.1365-2966.2009.15979.x. S2CID 119096173.
\item 格雷，R. O.；科巴利，C. J.；加里森，R. F.；麦克法登，M. T.；罗宾逊，P. E.（2003年10月）。“近邻恒星（NStars）项目贡献：40秒差距内M0型以前恒星的光谱观测：北部样本。I”。天文学杂志. 126 (4): 2048–2059. arXiv:astro-ph/0308182. Bibcode:2003AJ....126.2048G. doi:10.1086/378365. S2CID 119417105.
\item 谢纳夫林，V. I.；塔拉诺娃，O. G.；纳吉普，A. E.（2011年1月）。“炽热恒星周圍塵埃包層的搜尋與研究”。天文学报告. 55 (1): 31–81. 文献标识码:2011ARep...55...31S. doi:10.1134/S1063772911010070. S2CID 122700080.
\item 塞纳罗，A. J.；佩莱蒂耶，R. F.；桑切斯-布拉兹奎斯，P.；塞拉姆，S. O.；托洛巴，E.；卡迪埃尔，N.；法尔孔-巴罗索，J.；戈尔加斯，J.；希门尼斯-比森特，J.；瓦兹德基斯，A.（2007年1月）。“中分辨率艾萨克·牛顿望远镜实测光谱库——II。”恒星大气参数. 皇家天文学会月报. 374 (2): 664–690. arXiv:astro-ph/0611618. Bibcode:2007MNRAS.374..664C. doi:10.1111/j.1365-2966.2006.11196.x. S2CID 119428437.
\item 西恩，爱德华·M.；霍尔伯格，J. B.；奥斯瓦尔特，特里·D.；麦库克，乔治·P.；瓦萨托尼克，理查德（2009年12月）。“距离太阳20秒差距以内的白矮星：运动学与统计学”。天文学杂志. 138 (6): 1681–1689. arXiv:0910.1288. 文献标识码:2009AJ....138.1681S. DOI:10.1088/0004-6256/138/6/1681. S2CID 119284418.
\item D.S. 海斯；L.E. 帕西内蒂；A.G. 戴维斯·菲利普（2012年12月6日）。恒星基本量的标定：1984年5月24日至29日于意大利科莫奥尔莫别墅举行的国际天文学联合会第111次研讨会论文集。施普林格科学与商业媒体。第129页及以后ISBN 978-94-009-5456-4.
\item 阿里亚斯，朱莉娅·I. 等人（2016年8月）。“银河系O型恒星光谱巡天（GOSSS）中OVz型恒星的光谱分类与性质”. 《天文期刊》. 152 (2): 31. arXiv:1604.03842. 文献标识码:2016AJ....152...31A. DOI:10.3847/0004-6256/152/2/31. S2CID 119259952.
\item 麦克罗伯特，艾伦（2006年8月1日）。“恒星的光谱类型”. 天空与望远镜.
\item 艾伦，J. S.“恒星光谱的分类”.伦敦大学学院物理与天文学系：天体物理组。检索日期：2014年1月1日
\item 马伊兹·阿佩亚尼兹，J.；沃尔本，诺兰·R.；莫雷尔，N. I.；尼梅拉，V. S.；尼兰，E. P.（2007）。“皮斯米斯24-1：恒星质量上限的保持”。天体物理学报. 660 (2): 1480–1485. arXiv:astro-ph/0612012. 文献标识码:2007ApJ...660.1480M. DOI:10.1086/513098. S2CID 15936535.
\item 沃尔本，诺兰·R.；索塔，阿尔弗雷多；迈斯·阿佩亚尼兹，赫苏斯；阿尔法罗，埃米利奥·J.；莫雷尔，妮迪娅·I.；巴尔巴，罗多尔福·H.；阿里亚斯，胡莉娅·I.；加门，罗伯托·C.（2010）。“银河系O型恒星光谱巡天的早期结果：光谱中的C III发射线”。天体物理学杂志快报. 711(2): L143.arXiv:1002.3293. Bibcode:2010ApJ...711L.143W. doi:10.1088/2041-8205/711/2/L143. S2CID 119122481.
\item 法里尼亚，塞西莉亚；博施，吉列尔莫·L.；莫雷尔，尼迪娅·I.；巴尔巴，罗多尔福·H.；沃尔本，诺兰·R.（2009）。“大麦哲伦云中N159/N160复合体的光谱研究”。天文学杂志. 138 (2): 510–516. arXiv:0907.1033. 文献标识码:2009AJ....138..510F. DOI:10.1088/0004-6256/138/2/510. S2CID 18844754.
\item 劳夫，G.；曼弗鲁瓦，J.；戈塞，E.；纳泽，Y.；萨纳，H.；德贝克，M.；福尔米，C.；莫法特，A. F. J.（2007）。“年轻疏散星团韦斯特伦德2核心中的早型恒星”。天文学与天体物理学. 463 (3): 981–991. arXiv:astro-ph/0612622. 文献标识码:2007A&A...463..981R. DOI:10.1051/0004-6361:20066495. S2CID 17776145.
\item 克劳瑟，保罗·A.（2007）。“沃尔夫-拉叶星的物理性质”。天文学与天体物理学年评. 45 (1): 177–219. arXiv:astro-ph/0610356. 文献标识码:2007ARA&A..45..177C. DOI:10.1146/annurev.astro.45.051806.110615. S2CID 1076292.
\item 朗特里·莱什，J.（1968）。“古尔德带的运动学：一个正在膨胀的星群？”. 天体物理学杂志增刊系列. 17: 371. 文献标识码:1968ApJS...17..371L. DOI:10.1086/190179.
\item 若干恒星光谱分析及太阳黑子新观测, P. 塞基,法国科学院会议记录 63（1866年7月至12月），第364–368页。
\item 关于恒星光谱分析的新研究, P. 塞基,法国科学院院刊 63（1866年7月至12月），第621–628页。
\item 赫恩肖，J. B.（1986）。星光的分析：一百五十载天文学光谱学。英国剑桥：剑桥大学出版社。第60页、第134页。ISBN 978-0-521-25548-6.
\item 恒星光谱的分类：一些历史.
\item 卡勒，詹姆斯·B.（1997）。恒星及其光谱：光谱序列导论。剑桥：剑桥大学出版社。第62–63. ISBN 978-0-521-58570-5.
\item 第60–63页，赫恩肖1986年；第623–625页，塞基1866年。
\item 第62–63页，赫恩肖 1986。
\item 第60页，赫恩肖 1986。
\item 捕捉光明：首批拍摄天空的男女们被遗忘的人生作者：斯特凡·休斯。
\item 皮克林，爱德华·C.（1890）。“作为亨利·德雷珀纪念的一部分，用8英寸巴奇望远镜拍摄的恒星光谱德雷珀目录”。哈佛学院天文台年鉴. 27: 1. 文献标识码:1890AnHar..27....1P.
\item 第106–108页，赫恩肖 1986年。
\item 佩恩，塞西莉亚·H.（1930）。“O型恒星的分类”。哈佛学院天文台公报. 878: 1. 文献标识码:1930BHarO.878....1P.
\item “威廉米娜·弗莱明”.牛津参考. 检索日期2020年6月10日
\item “威廉米娜·帕顿·弗莱明 -”.www.projectcontinua.org. 检索日期2020年6月10日
\item “恒星光谱分类”.spiff.rit.edu。检索日期：2020年6月10日
\item 赫恩肖（1986年）第111–112页
\item 莫里，安东尼娅·C.；皮克林，爱德华·C.（1897）。“作为亨利·德雷珀纪念项目的一部分，用11英寸德雷珀望远镜拍摄的明亮恒星光谱”。哈佛学院天文台年鉴. 28: 1. Bibcode:1897AnHar..28....1M.
\item “安东尼娅·莫里”.www.projectcontinua.org。检索日期：2020年6月10日
赫恩肖，J.B.（2014年3月17日）。恒星光谱分析：两个世纪的天文学光谱学（第2版）。纽约，纽约州。ISBN 978-1-107-03174-6. OCLC 855909920.
格雷，理查德·O.；科巴利，克里斯托弗·J.；伯加瑟，亚当·J.（2009）。恒星光谱分类。新泽西州普林斯顿：普林斯顿大学出版社。ISBN 978-0-691-12510-7. OCLC 276340686.
\item 琼斯，贝茜·扎班；博伊德，莱尔·吉福德（1971）。哈佛学院天文台：前四任台长，1839-1919（第1版）。剑桥：哈佛大学出版社M.A.贝尔克纳普出版社。ISBN 978-0-674-41880-6. OCLC 1013948519.
\item 卡农，安妮·J.；皮克林，爱德华·C.（1901）。“作为亨利·德雷珀纪念项目的一部分，用13英寸博伊登望远镜拍摄的明亮南天恒星光谱”。哈佛学院天文台年鉴. 28: 129. 文献标识码:1901AnHar..28..129C.
\item 赫恩肖（1986）第117–119页，
\item 卡农，安妮·贾普；皮克林，爱德华·查尔斯（1912）。“利用光谱对1,688颗南天恒星进行分类”。哈佛学院天文台年鉴. 56 (5): 115. 文献标识码:1912AnHar..56..115C.
\item 赫恩肖（1986）第121–122页
\item “安妮·贾普·坎农”.www.projectcontinua.org。检索日期：2020年6月10日
\item 纳索，J. J.；赛弗特，卡尔·K.（1946年3月）。“北极五度范围内BD型恒星的光谱”。天体物理学报. 103: 117. 文献标识码:1946ApJ...103..117N. doi:10.1086/144796.
\item 菲茨杰拉德，M. 皮姆（1969年10月）。“威尔逊山与摩根－基南分类系统中光谱-光度等级的比较”。加拿大皇家天文学会期刊. 63: 251. 文献标识码:1969JRASC..63..251P.
\item 桑达奇，A.（1969年12月）。“新型亚矮星。II. 具有大自行的112颗恒星的径向速度、测光及初步空间运动”. 天体物理学报. 158: 1115. 文献标识码:1969ApJ...158.1115S. DOI:10.1086/150271.
\item 诺里斯，杰克逊·M.；赖特，贾森·T.；韦德，理查德·A.；马哈德万，苏夫拉特；盖特尔，萨拉（2011年12月）。“未探测到HD 149382的疑似亚恒星伴星”。天体物理学杂志. 743 (1). 88. arXiv:1110.1384. Bibcode:2011ApJ...743...88N. doi:10.1088/0004-637X/743/1/88. S2CID 118337277.
\item 加里森，R. F. (1994). “MK流程的标准层级体系” （PDF）。载于科巴利，C. J.；格雷，R. O.；加里森，R. F.（编）。MK流程50年：一种强大的天体物理洞察工具。太平洋天文学会会议系列。第60卷。旧金山：太平洋天文学会。第3–14. ISBN 978-1-58381-396-6. OCLC 680222523.
\item 亲爱的，大卫。“晚期恒星”。科学互联网百科。检索日期：2007年10月14日
\item 沃尔本，N. R.（2008）。“OB光谱的多波长系统学”。大质量恒星：基本参数与恒星周物质相互作用（编者：P. 贝纳利亚. 33: 5. 文献标识码:2008RMxAC..33....5W.
\item 恒星光谱图集，附光谱分类概要, W. W. 摩根、P. C. 基南和E. 凯尔曼，芝加哥：芝加哥大学出版社，1943年。
\item 沃尔本，N. R.（1971）。“OB型恒星的一些光谱特征：对某些OB型恒星空间分布及分类参考系的探究”. 天体物理学杂志增刊系列. 23: 257. 文献编码:1971ApJS...23..257W. DOI:10.1086/190239.
\item 摩根，W. W.；阿布特，赫尔穆特 A.；塔普斯科特，J. W.（1978）。“太阳之前恒星的修订MK光谱图集”。威廉姆斯湾：耶克斯天文台. 文献标识码:1978rmsa.book.....M.
\item 沃尔本，诺兰·R.；霍沃思，伊恩·D.；莱农，丹尼尔·J.；马西，菲利普；欧伊，M. S.；莫法特，安东尼·F. J.；斯卡尔科夫斯基，格温；莫雷尔，妮迪娅·I.；德里森，洛朗；帕克，乔尔·Wm.（2002）。“针对最早期O型恒星的新光谱分类系统：O2型的定义” （PDF）. 天文学杂志. 123 (5): 2754–2771. 文献标识码:2002AJ....123.2754W. DOI:10.1086/339831. S2CID 122127697.
\item 伊丽莎白·豪威尔（2013年9月21日）。“角宿一：王者之星”。Space.com。检索日期：2022年4月13日
\item 斯莱特巴克，阿尔内（1988年7月）。“贝星”. 太平洋天文学会出版物. 100: 770–784. 文献标识码:1988PASP..100..770S. DOI:10.1086/132234.
\item “距离最近的100个恒星系统”.www.astro.gsu.edu。检索日期：2022年4月13日
\item “20光年以内的恒星”.
\item 摩根，W. W.；基南，P. C.（1973）。“光谱分类”。天文学与天体物理学年评. 11: 29. 文献标识码:1973ARA&A..11...29M. DOI:10.1146/annurev.aa.11.090173.000333.
\item 摩根，W. W.；阿布特，赫尔穆特 A.；塔普斯科特，J. W.（1978）。修订版MK恒星光谱图集：适用于比太阳更早的恒星。芝加哥大学耶克斯天文台。文献标识码:1978rmsa.book.....M.
\item 格雷，R. O；加里森，R. F（1989）。“早期F型恒星——精化分类、与斯特龙格伦测光法的对比以及自转效应”。天体物理学杂志增刊系列. 69: 301. 文献标识码:1989ApJS...69..301G. DOI:10.1086/191315.
\item 基南，菲利普·C.；麦克尼尔，雷蒙德·C.（1989）。“针对较冷恒星的修订MK型别珀金斯目录”。天体物理学杂志增刊系列. 71: 245. 文献标识码:1989ApJS...71..245K. DOI:10.1086/191373. S2CID 123149047.
\item 尼尤文海森，H.；德雅格尔，C.（2000）。“检验黄色演化空缺。三颗演化临界超巨星：HD 33579、HR 8752与IRC +10420”。天文学与天体物理学. 353: 163. Bibcode:2000A&A...353..163N.
\item “从宇宙时间尺度来看，地球的宜居期已接近尾声 | 国际太空 fellowship”. Spacefellowship.com. 检索日期：2012年5月22日.
\item “金发姑娘”恒星或许正适合寻找宜居世界。NASA.com，2019年3月7日。检索日期：2021年8月26日。
\item “发现：恒星与人体一样凉爽 | 科学任务理事会”.science.nasa.gov。于2011年10月7日从原始网页存档。检索日期：2017年7月12日
\item “银河翻新”.www.spacetelescope.org. 欧洲航天局/哈勃. 检索日期：2015年4月29日
\item 佩恩，塞西莉亚·H.（1930）。“O型恒星的分类”。哈佛学院天文台公报. 878: 1. 文献标识码:1930BHarO.878....1P.
\item 菲格，唐纳德·F.；麦克莱恩，伊恩·S.；纳哈罗，弗朗西斯科（1997）。“沃尔夫-拉叶星AK波段光谱图集”. 天体物理学报. 486 (1): 420–434. 文献标识码:1997ApJ...486..420F. DOI:10.1086/304488.
\item 金斯伯格，R. L.；巴洛，M. J.；斯托里，P. J.（1995）。“WO沃尔夫-拉叶星的性质”。天文学与天体物理学. 295: 75. 文献标识码:1995A&A...295...75K.
\item 廷克勒，C. M.；拉默斯，H. J. G. L. M.（2002）。“富含氢的行星状星云中心恒星的质量损失率能否用作距离指示器？”. 天文学与天体物理学. 384 (3): 987–998. 文献标识码:2002A&A...384..987T. DOI:10.1051/0004-6361:20020061.
\item 米扎尔斯基，B.；克劳瑟，P. A.；德马可，O.；科彭，J.；莫法特，A. F. J.；阿克尔，A.；希尔维格，T. C.（2012）。“IC 4663：首个明确的[WN]沃尔夫-拉叶型行星状星云中心恒星”. 皇家天文学会月报. 423 (1): 934–947. arXiv:1203.3303. 文献标识码:2012MNRAS.423..934M. DOI:10.1111/j.1365-2966.2012.20929.x. S2CID 10264296.
\item 克劳瑟，P. A.；沃尔本，N. R.（2011）。“O2-3.5 If*/WN5-7型恒星的光谱分类”。《皇家天文学会月报》416（2）：1311–1323。arXiv：1105.4757文献标识码2011MNRAS.416.1311CDOI10.1111/j.1365-2966.2011.19129.xS2CID118455138
\item 柯克帕特里克，J. D.（2008）。“我们对L型、T型和Y型矮星理解中的突出问题”。第14届剑桥低温恒星研讨会. 384: 85. arXiv:0704.1522. 天文学文献标识码:2008ASPC..384...85K.
\item 柯克帕特里克，J·戴维；里德，I·尼尔；利伯特，詹姆斯；库特里，罗克·M.；纳尔逊，布兰特；贝希曼，查尔斯·A.；达恩，科纳德·C.；莫内，大卫·G.；吉齐斯，约翰·E.；斯克鲁茨基，迈克尔·F.（1999年7月10日）。“比M型更冷的矮星：利用2微米全天巡天（2MASS）发现定义光谱类型L”. 天体物理学报. 519 (2): 802–833. 文献标识码:1999ApJ...519..802K. DOI:10.1086/307414.
\item 柯克帕特里克，J·戴维（2005）。“新型光谱类型L和T” （PDF）. 天文学与天体物理学年评. 43 (1): 195–246. 文献标识码:2005ARA&A..43..195K. DOI:10.1146/annurev.astro.42.053102.134017. S2CID 122318616.
\item 柯克帕特里克，J·戴维；巴曼，特拉维斯·S；伯格瑟，亚当·J；麦戈文，马克·R；麦克莱恩，伊恩·S；廷尼，克里斯托弗·G；洛厄朗斯，帕特里克·J（2006）。“发现一颗极年轻的场L型矮星，2MASS J01415823−4633574”。天体物理学杂志. 639 (2): 1120–1128. arXiv:astro-ph/0511462. Bibcode:2006ApJ...639.1120K. doi:10.1086/499622. S2CID 13075577.
 卡门辛德，马克斯（2006年9月27日）。“恒星光谱的分类及其物理解释” （PDF）. 天体实验室国王椅州立天文台：6 – 经海德堡大学提供。
 柯克帕特里克，J·戴维；库欣，迈克尔·C；杰利诺，克里斯托弗·R；贝希曼，查尔斯·A；廷尼，C·G；法瑟蒂，贾奎琳·K；施耐德，亚当；梅斯，格雷戈里·N（2013）。“发现Y1型矮星WISE J064723.23-623235.5”。天体物理学杂志. 776 (2): 128. arXiv:1308.5372. Bibcode:2013ApJ...776..128K. doi:10.1088/0004-637X/776/2/128. S2CID 6230841.
 迪肯，N. R.；汉布利，N. C.（2006）。“超冷矮星的Y光谱分类”。皇家天文学会月报371：1722–1730。arXivastro-ph/0607305doi10.1111/j.1365-2966.2006.10795.xS2CID14081778
 韦纳，迈克（2011年8月24日）。“美国国家航空航天局发现，一些恒星的温度比人体还低 | 技术新闻博客——雅虎！加拿大新闻”。Ca.news.yahoo.com。检索日期：2012年5月22日。
 文顿，丹妮尔（2011年8月23日）。“美国宇航局卫星发现迄今最冷、最暗的恒星”. 连线——经由www.wired.com。
 “美国国家航空航天局——NASA的WISE任务发现最冷的一类恒星”.www.nasa.gov。于2021年2月14日存档自原文。检索日期：2019年11月1日
 祖克曼，B.；宋，I.（2009）。“最小简氏质量、棕矮星伴星初始质量函数以及Y型矮星探测的预测”。天文学与天体物理学. 493 (3): 1149–1154. arXiv:0811.0429. 文献编码:2009A&A...493.1149Z. doi:10.1051/0004-6361:200810038. S2CID 18147550.
 杜普伊，T. J.；克劳斯，A. L.（2013）。“已知最冷次恒星天体的距离、光度与温度”。科学. 341 (6153): 1492–5. arXiv:1309.1422. Bibcode:2013Sci...341.1492D. doi:10.1126/science.1241917. PMID 24009359. S2CID 30379513.
 莱格特，桑迪·K.；库欣，迈克尔·C.；索蒙，迪迪埃；马利，马克·S.；罗埃利格，托马斯·L.；沃伦，斯蒂芬·J.；伯宁厄姆，本；琼斯，休·R·A.；柯克帕特里克，J·戴维；洛迪厄，尼古拉斯；卢卡斯，菲利普·W.；迈因泽，艾米·K.；马丁，爱德华多·L.；麦考格伦，马克·J.；平菲尔德，大卫·J.；斯隆，格雷戈里·C.；斯马特，理查德·L.；田村元秀；范克利夫，杰弗里·E.（2009）。“四颗约600 K T型矮星的物理性质”。天体物理学杂志. 695 (2): 1517–1526. arXiv:0901.4093. 文献标识码:2009ApJ...695.1517L. doi:10.1088/0004-637X/695/2/1517. S2CID 44050900.
 德洛姆，菲利普；德尔福斯，哈维耶；阿尔贝，洛伊克；阿尔蒂戈，埃蒂安；福尔韦耶，蒂埃里；雷莱，塞琳；阿拉尔，弗朗斯；霍梅尔，德里克；罗宾，安妮·C.；威洛特，克里斯·J.；刘，迈克尔·C.；杜皮，特伦特·J.（2008）。“CFBDS J005910.90-011401.3：正迈向T-Y型棕矮星过渡？”天文学与天体物理学. 482 (3): 961–971. arXiv:0802.4387. Bibcode:2008A&A...482..961D. doi:10.1051/0004-6361:20079317. S2CID 847552.
 伯宁厄姆，本；平菲尔德，D. J.；莱格特，S. K.；田村，M.；卢卡斯，P. W.；霍迈尔，D.；戴-琼斯，A.；琼斯，H. R. A.；克拉克，J. R. A.；石井，M.；久住原，M.；洛迪厄，N.；萨帕特罗-奥索里奥，玛丽亚·罗莎；韦内曼斯，B. P.；莫特洛克，D. J.；巴拉多-伊-纳瓦斯奎斯，D.；马丁，爱德华多·L.；马加祖，安东尼奥（2008）。“探索低至约550 K的次恒星温度范围”. 皇家天文学会月报. 391 (1): 320–333. arXiv:0806.0067. Bibcode:2008MNRAS.391..320B. doi:10.1111/j.1365-2966.2008.13885.x. S2CID 1438322.
 欧洲南方天文台. “一对非常酷的棕矮星”, 2011年3月23日
 卢曼，凯文·L.；埃斯普林，塔兰·L.（2016年5月）。“已知最冷棕矮星的光谱能量分布”. 《天文学杂志》. 152 (3): 78. arXiv:1605.06655. 文献标识码:2016AJ....152...78L. doi:10.3847/0004-6256/152/3/78. S2CID 118577918.
 “光谱类型代码”.simbad.u-strasbg.fr。检索日期：2020年3月6日
 伯宁厄姆，本；史密斯，L.；卡多索，C.V.；卢卡斯，P.W.；伯格瑟，亚当·J.；琼斯，H.R.A.；斯马特，R.L.（2014年5月）。“一颗T6.5亚矮星的发现”. 皇家天文学会月报. 440 (1): 359–364. arXiv:1401.5982. 文献标识码:2014MNRAS.440..359B. DOI:10.1093/mnras/stu184. ISSN 0035-8711. S2CID 119283917.
 克鲁兹，凯勒·L.；柯克帕特里克，J·戴维；伯格瑟，亚当·J.（2009年2月）。“在野外识别出的年轻L型矮星：从L0到L5的初步低重力光学光谱序列”。天文学杂志. 137 (2): 3345–3357. arXiv:0812.0364. 文献标识码:2009AJ....137.3345C. DOI:10.1088/0004-6256/137/2/3345. ISSN 0004-6256. S2CID 15376964.
 洛珀，达格妮·L.；柯克帕特里克，J·戴维；库特里，罗克·M.；巴尔曼，特拉维斯；伯加瑟，亚当·J.；库欣，迈克尔·C.；罗埃利格，托马斯；麦戈文，马克·R.；麦克莱恩，伊恩·S.；赖斯，艾米丽；斯威夫特，布兰登·J.（2008年10月）：“从2MASS自行运动巡天中发现两颗近距离奇特的L型矮星：年轻还是富含金属？”天体物理学报. 686 (1): 528–541. arXiv:0806.1059. Bibcode:2008ApJ...686..528L. doi:10.1086/591025. ISSN 0004-637X. S2CID 18381182.
 柯克帕特里克，J·戴维；卢珀，达格妮·L；布尔加瑟，亚当·J；舒尔，史蒂文·D；库特里，罗克·M；卡什инг，迈克尔·C；克鲁兹，凯勒·L；斯威特，安妮·C；纳普，吉莉安·R；巴尔曼，特拉维斯·S；博昌斯基，约翰·J（2010年9月）。“利用多 epoch 两微米全天巡天数据开展近红外自行测量巡天的发现”。天体物理学杂志增刊系列. 190 (1): 100–146. arXiv:1008.3591. Bibcode:2010ApJS..190..100K. doi:10.1088/0067-0049/190/1/100. ISSN 0067-0049. S2CID 118435904.
 法赫蒂，杰奎琳·K.；里德尔，阿德里克·R.；克鲁兹，凯勒·L.；加涅，乔纳森；菲利帕佐，约瑟夫·C.；兰布里德斯，埃里尼；菲卡，哈莉；温伯格，艾莉西亚；索尔斯滕森，约翰·R.；廷尼，C.G.；巴尔达萨雷，维维安（2016年7月）。“系外行星类褐矮星的种群特性”. 天体物理学杂志增刊系列. 225 (1): 10. arXiv:1605.07927. 文献标识码:2016ApJS..225...10F. DOI:10.3847/0067-0049/225/1/10. ISSN 0067-0049. S2CID 118446190.
 “颜色-星等数据”.太空望远镜科学研究所(www.stsci.edu)。检索日期：2020年3月6日
 布伊格，R.（1954）。天体物理学年报，第17卷，第104页
 基南，P. C.（1954）。“S型恒星的分类”。天体物理学报. 120: 484. 文献标识码:1954ApJ...120..484K. DOI:10.1086/145937.
 西恩，E. M.；格林斯坦，J. L.；兰德斯特里特，J. D.；利伯特，詹姆斯；希普曼，H. L.；韦格纳，G. A.（1983）。“一种提议的新型白矮星光谱分类系统”. 天体物理学报. 269: 253. 文献标识码:1983ApJ...269..253S. doi:10.1086/161036.
 科尔西科，A. H.；阿尔陶斯，L. G.（2004）。“脉动DB型白矮星的周期变化速率”。天文学与天体物理学. 428: 159–170. arXiv:astro-ph/0408237. Bibcode:2004A&A...428..159C. doi:10.1051/0004-6361:20041372. S2CID 14653913.
 麦库克，乔治·P.；锡恩，爱德华·M.（1999）。“光谱已确认的白矮星目录”。天体物理学杂志增刊系列. 121 (1): 1–130. 文献标识码:1999ApJS..121....1M. CiteSeerX 10.1.1.565.5507. doi:10.1086/313186. S2CID 122286998.
 阿佩利亚尼兹，J. 马伊斯；巴尔瓦，R. H.；阿兰达，R. 费尔南德斯；冈萨雷斯，M. 潘塔莱奥尼；贝利多，P. 克雷斯波；索塔，A.；阿尔法罗，E. J.（2022年1月1日）。“维拉弗兰卡银河OB星团目录——II：从盖亚DR2到EDR3及十个含O型恒星的新系统”. 天文学与天体物理学. 657: A131.arXiv:2110.01464. 文献标识码:2022A&A...657A.131M. DOI:10.1051/0004-6361/202142364. ISSN 0004-6361.
 马西，菲利普；纽根特，凯瑟琳·F.；斯马特，布里安娜·M.（2016年9月1日）。“对M31和M33中大质量恒星的光谱巡天研究*”. 天文期刊. 152 (3): 62. arXiv:1604.00112. 文献标识码:2016AJ....152...62M. DOI:10.3847/0004-6256/152/3/62. ISSN 0004-6256.
 “脉动变星与赫罗图（H-R图）”。哈佛-史密森天体物理中心，2015年3月9日。检索日期：2016年7月23日。
 雅科夫列夫，D. G.；卡明克尔，A. D.；哈恩塞尔，P.；格涅金，O. Y. （2002）。“3C 58中冷却中的中子星”。天文学与天体物理学. 389: L24 – L27. arXiv:astro-ph/0204233. Bibcode:2002A&A...389L..24Y. doi:10.1051/0004-6361:20020699. S2CID 6247160.
 “恒星与宜居行星”. www.solstation.com.
\end{enumerate}
\subsection{延伸阅读}
\begin{itemize}
\item 哈雷，扬-文森特；赫勒，勒内（2021）。“恒星的数字颜色编码”。天文消息. 342 (3): 578–587. arXiv:2101.06254. Bibcode:2021AN....342..578H. doi:10.1002/asna.202113868. S2CID 231627588.
\end{itemize}